\section{Exact inversion from four fixed physical windows}
\label{sec:v8-fixed}
We now keep the offsets fixed while the requested accuracy improves.
This experiment uses the exact analytic support family
\eqref{eq:g-isogap-family}, not an unspecified smooth remainder in its
onset expansion. The true gap is unknown. The reference circle has
$R_c=1/4$ and $g_c=1/2$. Put $x_i=r_i-R$, and let $e=(e_1,e_2,e_3)$
be the elementary symmetric functions of the three radius increments.
Thus $R=(1-g)/2$, the unordered curvatures are $(R+x_i)^{-1}$, and
$A=A(R,e)$ is the exact function in \eqref{eq:v6-area-constraint}.
Throughout this section the physical parameter set is a sufficiently
small fixed compact neighborhood of the reference member. In particular
it contains coalescing and repeated-root configurations. It is not a
set of labelled root parameters with a separation condition.

\subsection{Analytic descent through coalescence}
For $j=1,2,3$, define at positive offset
\[
 F_j(g,x,d)=d^{-2}\Pr_{g,x}\{N_{jg+d}\ge j+1\}.
\]
These are normalized exact finite-offset probabilities. The analytic
extension needed for inversion at a triple root is stronger than
continuity in the labelled radii.

\begin{lemma}[Symmetric analytic finite-offset coordinates]
\label{lem:v8-descent}
The functions $F_j$ have jointly analytic extensions to a neighborhood
of $(g_c,0,0)$ in the variables $(g,e,d)$, before restriction to the
real-rooted physical locus. On a smaller such neighborhood,
\begin{equation}\label{eq:v8-analytic-expansion}
                     F_j(g,e,d)=C_j(g,e)+d B_j(g,e,d),
\end{equation}
where $B_j$ is analytic and $C_j$ is the physical leading-amplitude map
in Theorem~\ref{thm:v6-three}. Every fixed derivative is bounded on a
smaller closed neighborhood.
\end{lemma}
\begin{proof}
The linear map from $(\alpha,\beta,\zeta)$ to $(x_1,x_2,x_3)$ is
invertible by \eqref{eq:g-contact-curvatures}. The support family is
analytic in these variables and $g$. For each of the finitely many
channels and flight numbers under consideration,
Lemma~\ref{lem:g-relative} gives the analytic stationary action and
its twist on a common complex neighborhood. The positive real endpoint
Hessian has an analytic square-root branch on a smaller complex
neighborhood. The constructive Morse map in Lemma~\ref{lem:g-radial}
and its inverse are consequently analytic there, by the analytic
inverse-function theorem.

After physical residual-time integration and the Morse change of
variables, the probability divided by $d^2$ is an integral over a fixed
disk of an analytic amplitude at $\sqrt{2d}\,z$, with weight
$1-|z|^2$. Expand the amplitude in its convergent endpoint power series
on a strictly smaller polydisk. The series is normally convergent,
uniformly in the remaining complex parameters. Integration annihilates
all odd total degrees and leaves a normally convergent series in integer
powers of $d$. This constructs an analytic extension across $d=0$;
it does not assign a nonzero physical probability below onset. The six
channels all have exactly the same gap in the present family, so adding
them introduces no moving activation interface.

A rotation of the triangular lattice by $\pi/3$ cycles the three radius
increments. Reflection in the horizontal axis fixes the first and
interchanges the other two. These transformations generate their full
permutation group. They preserve phase volume and the unlabelled count
event, hence $F_j(g,x,d)$ is symmetric in $x$ on the real physical open
set. The analytic identity principle gives the same symmetry on a
permutation-invariant complex polydisk.

We spell out descent at the discriminant. Choose a smaller coefficient
polydisk on which all roots of
\[
                     z^3-e_1z^2+e_2z-e_3
\]
stay strictly inside the original root polydisk; this is possible by
continuity of the root multiset at $e=0$. Away from the discriminant,
local holomorphic root branches define $F_j$ as a holomorphic function
of $(g,e,d)$. Symmetry makes this definition independent of the branches
and their permutations. It is locally bounded uniformly up to the
discriminant, since all roots stay in the smaller root polydisk and the
original function is bounded on its closure. The removable-singularity
theorem for locally bounded holomorphic functions across an analytic
hypersurface therefore extends it jointly across the discriminant.
One can apply the theorem in the product domain with $g,d$ as additional
variables, so joint analyticity is retained.

At $d=0$, the physical onset formula identifies the extension with
$C_j(g,e)$ on distinct real-rooted configurations. Such configurations
contain real open subsets of the coefficient domain, so uniqueness of
analytic extension identifies the two analytic functions everywhere
near $(g_c,0)$. Dividing their difference by $d$, or using its convergent
Taylor series, gives \eqref{eq:v8-analytic-expansion}. Cauchy estimates
on smaller polydisks give the final boundedness assertion.
\end{proof}

\subsection{The unknown gap as a fourth coordinate}
Fix a positive number $a$, to be chosen sufficiently small once and
then kept independent of accuracy and confidence. Use the four windows
\begin{equation}\label{eq:v8-windows}
 (j,t)=(1,g_c+a),\quad(1,g_c+2a),\quad
                         (2,2g_c+a),\quad(3,3g_c+a).
\end{equation}
Their normalizations use the known programmed excesses, not the unknown
actual ones. If $v=g-g_c$, put
\begin{equation}\label{eq:v8-Q}
 Q_a(g,e)=\begin{pmatrix}
 (a-v)^2F_1(g,e,a-v)/a^2\\
 (2a-v)^2F_1(g,e,2a-v)/(4a^2)\\
 (a-2v)^2F_2(g,e,a-2v)/a^2\\
 (a-3v)^2F_3(g,e,a-3v)/a^2
 \end{pmatrix}.
\end{equation}
On $|v|<a/6$ these are the normalized physical means of the four
binary reports; all actual offsets are positive. Decrease $a$ and the
neighborhood so they lie in the common physical collar.

\begin{theorem}[Four-window physical inverse]\label{thm:v8-four}
For some fixed $a>0$, the exact map $Q_a$ is analytically invertible
on a full neighborhood of $(g_c,0)$ in coefficient coordinates. There
is a compact physical neighborhood $K$ of the circular member such
that, for $\theta,\widetilde\theta\in K$,
\begin{align}
 |g-\widetilde g|+\|e-\widetilde e\|_\infty+|A-\widetilde A|
       &\le C\|Q_a(\theta)-Q_a(\widetilde\theta)\|_\infty,
                                        \label{eq:v8-fixed-stability}\\
 \operatorname{dist}_{\rm match}(\kappa,\widetilde\kappa)
       &\le C\|Q_a(\theta)-Q_a(\widetilde\theta)\|_\infty^{1/3}.
                                        \label{eq:v8-fixed-roots}
\end{align}
The neighborhood and constants do not depend on a root separation.
Only the known family formula, reference window constants, and the
fixed neighborhood are supplied. The true gap, area, channel labels,
and impact positions are not supplied.
\end{theorem}
\begin{proof}
Let $D=D_e(C_1,C_2,C_3)(g_c,0)$. By
\eqref{eq:v6-three-determinant},
\begin{equation}\label{eq:v8-D}
 \det D=-\frac{2\sqrt2(15804720 A_c+64253\pi)}{72930375 A_c^4},
       \qquad A_c=\frac{\sqrt3}{2}-\frac\pi{16}>0.
\end{equation}
In particular $\det D\ne0$. The derivation of this exact determinant,
including its twelve amplitude derivatives and the physical area
constraint, is given in Theorem~\ref{thm:v6-three}.

For a window of order $j$ with programmed excess $la$ at the reference
gap, differentiation at $(g_c,0)$ gives
\[
 \begin{aligned}
 \partial_g\left[\frac{(la-jv)^2}{(la)^2}
                  F_j(g,e,la-jv)\right]
       &=-\frac{2jC_j}{la}+O(1),\\
 D_e\left[\frac{(la-jv)^2}{(la)^2}
                  F_j(g,e,la-jv)\right]&=D_eC_j+O(a).
 \end{aligned}
\]
Here Lemma~\ref{lem:v8-descent} bounds derivatives in the full
coefficient neighborhood; estimates only in the labelled roots would
not justify the second formula at coalescence.

In $DQ_a$, subtract row one from row two, place this difference first,
and multiply the gap column by $a$. The resulting matrix converges as
$a\downarrow0$ to
\begin{equation}\label{eq:v8-block}
 \begin{pmatrix}
 C_1&0&0&0\\
 -2C_1&D_{11}&D_{12}&D_{13}\\
 -4C_2&D_{21}&D_{22}&D_{23}\\
 -6C_3&D_{31}&D_{32}&D_{33}
 \end{pmatrix}_{(g_c,0)}.
\end{equation}
Its determinant is $C_1\det D\ne0$. These row and column operations
are invertible for every $a>0$. Fix $a$ small enough that $DQ_a$ is
nonsingular. The analytic inverse-function theorem now applies to the
exact finite-offset map, not to its leading approximation.

Choose a convex data ball with closure inside the image of an inverse
chart, and a smaller compact physical neighborhood $K$ whose image
lies in that ball. The inverse has bounded derivative on the ball,
giving the gap and coefficient bounds in
\eqref{eq:v8-fixed-stability}. The analytic area identity gives its area
bound. The root-matching argument of Theorem~\ref{thm:v5-pairwise}
gives $\operatorname{dist}_{\rm match}(x,\widetilde x)
\le C\|e-\widetilde e\|_\infty^{1/3}$. Adding the change in
$R=(1-g)/2$ and using the bounded derivative of $r\mapsto r^{-1}$ on
the common positive radius interval proves \eqref{eq:v8-fixed-roots}.
Both the analytic inverse and the root estimate were established before
any exclusion of a repeated root, so no such exclusion is needed.
\end{proof}

\begin{corollary}[Deterministic-budget fixed-window acquisition]
\label{cor:v8-fixed-sampling}
There are constants $c,C>0$ such that for $0<\varepsilon,\delta<c$
and $0<\eta<1/4$, four-window binary observations give a Borel estimator
satisfying, uniformly on $K$,
\begin{equation}\label{eq:v8-joint-confidence}
 \Pr\{\operatorname{dist}_{\rm match}(\widehat\kappa,\kappa)
       \le\varepsilon,\quad |\widehat g-g|\le\delta\}\ge1-\eta,
\end{equation}
using a deterministic number of independent preparations bounded by
\begin{equation}\label{eq:v8-fixed-cost}
                  C(\varepsilon^{-6}+\delta^{-2})\log(C/\eta).
\end{equation}
The same estimator has coefficient and area errors at most
$C\min\{\varepsilon^3,\delta\}$ on this event. In particular,
curvature-only accuracy $\varepsilon$ costs at most
$C\varepsilon^{-6}\log(C/\eta)$ and also gives gap and area errors
$O(\varepsilon^3)$.
\end{corollary}
\begin{proof}
Take $n$ fresh preparations in each window. The centered Bernoulli
moment-generating function is at most $\exp(t^2/8)$: the second
derivative of its log is the variance of a tilted Bernoulli law and
is at most $1/4$, while its value and first derivative at zero vanish.
Exponential Markov and optimization therefore give
$\Pr\{|\widehat p-p|>u\}\le2e^{-2nu^2}$. Dividing the four fractions
by their fixed design normalizers and taking a union bound gives
\[
 \Pr\{\|\widehat Q-Q_a(\theta)\|_\infty>
                       C_a\sqrt{\log(8/\eta)/n}\}\le\eta.
\]
The dependence on $a$ is a constant because $a$ is fixed.

Minimize $\|Q_a(\vartheta)-\widehat Q\|_\infty$ on $K$, using the
Borel selection of Lemma~\ref{lem:v7-selection}. This is defined on
all outcomes, including data outside the inverse chart. Since the
truth lies in $K$, its fitted exact data differ from the true data by
at most $2\|\widehat Q-Q_a(\theta)\|_\infty$. Apply
Theorem~\ref{thm:v8-four}. Taking
\[
 n=\left\lceil C\min\{\varepsilon^3,\delta\}^{-2}
                                     \log(C/\eta)\right\rceil
\]
with a sufficiently large constant proves all stated error bounds.
Finally $\min\{\varepsilon^3,\delta\}^{-2}
=\max\{\varepsilon^{-6},\delta^{-2}\}
\le\varepsilon^{-6}+\delta^{-2}$, and there are exactly $4n$
preparations. This specifies a statistical estimator, not an efficient
algorithm for exact evaluation of the nonlinear probability map.
\end{proof}

The fixed-offset analytic descent and the four-window design were
suggested in the referee comparison supplied with the two independent
reports of September 9, 2026. The preceding proofs incorporate that
comparison. The joint lower bound below identifies the additional
cost of demanding a finer gap estimate.
